\documentclass[a4paper,12pt]{article}
\usepackage[utf8]{inputenc}
\usepackage[T1]{fontenc}
\usepackage[ngerman]{babel}
\usepackage{amsmath,amssymb}
\usepackage{booktabs}
\usepackage{enumitem}
\usepackage[margin=2.5cm]{geometry}
\setlength{\emergencystretch}{2em}

\begin{document}

\begin{center}
{\LARGE\bfseries \"Ubungsblatt 6 -- L\"osungen}\\[0.5em]
{\large VC Logik}
\end{center}

\bigskip

%% ============================================================
%% AUFGABE 1
%% ============================================================
\subsection*{Aufgabe 1 \textnormal{(2 Punkte)} -- Doppelte Resolution ist inkorrekt}

\textbf{Was ist eine Klausel?}
Eine Klausel ist eine Oder-Verkn\"upfung von Literalen, also z.\,B.\
$X \vee \neg Y$. In Mengenschreibweise wird sie als $\{X, \neg Y\}$ notiert;
das Komma zwischen den Literalen bedeutet hier $\vee$. Die \textbf{leere
Klausel} $\emptyset$ ist eine Klausel ohne Literale -- sie ist immer falsch,
weil ein ODER ohne Bestandteile nichts hat, was sie wahr machen k\"onnte.

\medskip
\textbf{Was ist die normale Resolutionsregel?}
Aus zwei Klauseln $D_1$ und $D_2$, in denen ein Literal $\ell$ einmal
positiv und einmal negativ vorkommt ($\ell \in D_1$, $\neg\ell \in D_2$),
darf eine neue Klausel gebildet werden:
\[
\text{Resolvent} \;=\; (D_1 \setminus \{\ell\}) \cup (D_2 \setminus \{\neg\ell\}).
\]
Aus $D_1 = \{X, A\}$ und $D_2 = \{\neg X, B\}$ wird also $\{A, B\}$.
Diese Regel ist korrekt, denn:
\begin{itemize}[nosep]
\item Ist $X = 1$, muss in $D_2$ schon $B = 1$ sein, damit $D_2$ wahr ist.
\item Ist $X = 0$, muss in $D_1$ schon $A = 1$ sein, damit $D_1$ wahr ist.
\end{itemize}
In jedem Fall ist also $A \vee B$ wahr -- der Resolvent stimmt immer, wenn
die Pr\"amissen stimmen.

\medskip
\textbf{Was ist die ``doppelte'' Resolutionsregel?}
Die in der Aufgabe beschriebene Regel will gleich \emph{zwei} Literale
auf einmal eliminieren. Voraussetzung: in $C_1$ stehen zwei Literale $a, b$
und in $C_2$ zwei Literale $c, d$, sodass $a \equiv \neg c$ und
$b \equiv \neg d$ (also beide Paare komplement\"ar). Dann liefert die
Regel:
\[
\text{Doppelresolvent} \;=\; (C_1 \setminus \{a, b\}) \cup (C_2 \setminus \{c, d\}).
\]

\medskip
\textbf{Was hei\ss t ``inkorrekt''?}
Eine Inferenzregel ist \emph{korrekt}, wenn aus den Pr\"amissen die
Konklusion logisch folgt -- also: jede Belegung, die die Pr\"amissen wahr
macht, macht auch die Konklusion wahr. \emph{Inkorrekt} hei\ss t, dass es
mindestens \emph{eine} Belegung gibt, unter der die Pr\"amissen wahr sind,
die Konklusion aber falsch. Genau ein solches Gegenbeispiel reicht aus.

\medskip
\textbf{Aufgabenstellung.}
Gegeben Klauseln $C_1$ und $C_2$ mit $\{a,b\} \subseteq C_1$,
$\{c,d\} \subseteq C_2$, $a \equiv \neg c$ und $b \equiv \neg d$
(also z.\,B.\ $\{X, \neg Y\} \subseteq C_1$ und $\{\neg X, Y\} \subseteq C_2$).
Der Doppelresolvent ist
$(C_1 \setminus \{a,b\}) \cup (C_2 \setminus \{c,d\})$.
Zu zeigen ist, dass diese Regel inkorrekt ist.

\medskip
\textbf{Wahl des Gegenbeispiels.}
Es ist sinnvoll, die kleinstm\"oglichen Klauseln zu w\"ahlen, die die
Voraussetzung erf\"ullen. Das sind genau die Klauseln, die nichts au\ss er
den vier Literalen $a, b, c, d$ enthalten:
\[
C_1 = \{X,\, \neg Y\}, \qquad C_2 = \{\neg X,\, Y\}.
\]
Mit der Zuordnung $a = X$, $b = \neg Y$, $c = \neg X$, $d = Y$ wird
gepr\"uft, ob die Voraussetzungen passen:
\begin{itemize}[nosep]
\item $a \equiv \neg c$: $X \equiv \neg(\neg X) \equiv X$. \checkmark
\item $b \equiv \neg d$: $\neg Y \equiv \neg Y$. \checkmark
\end{itemize}
Die Voraussetzungen sind also erf\"ullt.

\medskip
\textbf{Schritt 1: Doppelresolventen ausrechnen.}
Da $\{a, b\}$ ganz $C_1$ ist und $\{c, d\}$ ganz $C_2$ ist, wird in beiden
Klauseln alles entfernt:
\[
\underbrace{(C_1 \setminus \{X, \neg Y\})}_{=\;\emptyset}
\;\cup\;
\underbrace{(C_2 \setminus \{\neg X, Y\})}_{=\;\emptyset}
\;=\; \emptyset.
\]
Der Doppelresolvent ist also die leere Klausel $\emptyset$. Diese ist
unter \emph{jeder} Belegung falsch.

\medskip
\textbf{Schritt 2: Erf\"ullbarkeit von $C_1 \wedge C_2$ pr\"ufen.}
W\"are die Regel korrekt, m\"usste aus $C_1 \wedge C_2$ logisch $\emptyset$
folgen. Da $\emptyset$ aber immer falsch ist, w\"urde das nur funktionieren,
wenn $C_1 \wedge C_2$ \emph{nie} wahr wird -- also unerf\"ullbar ist. Eine
Wahrheitstabelle zeigt das Gegenteil:

\begin{center}
\renewcommand{\arraystretch}{1.2}
\begin{tabular}{cc|cc|c|c}
\toprule
$X$ & $Y$ & $C_1 = X \vee \neg Y$ & $C_2 = \neg X \vee Y$ & $C_1 \wedge C_2$ & Doppelres. $\emptyset$ \\
\midrule
0 & 0 & $0 \vee 1 = 1$ & $1 \vee 0 = 1$ & $1$ & $0$ \\
0 & 1 & $0 \vee 0 = 0$ & $1 \vee 1 = 1$ & $0$ & $0$ \\
1 & 0 & $1 \vee 1 = 1$ & $0 \vee 0 = 0$ & $0$ & $0$ \\
1 & 1 & $1 \vee 0 = 1$ & $0 \vee 1 = 1$ & $1$ & $0$ \\
\bottomrule
\end{tabular}
\end{center}

In den Zeilen $X = Y = 0$ und $X = Y = 1$ ist $C_1 \wedge C_2$ wahr. Der
Doppelresolvent ist aber in jeder Zeile falsch (weil $\emptyset$ konstant
falsch ist). Es gibt also Belegungen, unter denen die Pr\"amissen wahr,
die Konklusion aber falsch ist.

\medskip
\textbf{Schluss.}
Damit folgt $\emptyset$ \emph{nicht} logisch aus $C_1 \wedge C_2$. Die
doppelte Resolutionsregel w\"urde aber genau das behaupten -- sie ist also
inkorrekt.

\medskip
\textbf{Was geht hier eigentlich schief?}
Die normale Resolution liefert f\"ur dieselben beiden Klauseln nur
Tautologien:
\begin{itemize}[nosep]
\item \"Uber $X$: $\{\neg Y\} \cup \{Y\} = \{\neg Y, Y\}$ -- Tautologie.
\item \"Uber $Y$: $\{X\} \cup \{\neg X\} = \{X, \neg X\}$ -- Tautologie.
\end{itemize}
Die normale Resolution erkennt also korrekt, dass aus $C_1, C_2$ nichts
Neues folgt. Die doppelte Resolution \"uberspringt diesen Schritt und
schmei\ss t einfach beide komplement\"aren Paare gleichzeitig raus -- und
landet so bei $\emptyset$, einer viel st\"arkeren (und hier falschen)
Behauptung.

\medskip
\noindent\fbox{\parbox{\dimexpr\textwidth-2\fboxsep-2\fboxrule}{%
\textbf{Ergebnis Aufgabe 1.}\\[2pt]
F\"ur $C_1 = \{X, \neg Y\}$ und $C_2 = \{\neg X, Y\}$ liefert die doppelte
Resolutionsregel die leere Klausel $\emptyset$. Die Belegung $X = 1, Y = 1$
(und auch $X = 0, Y = 0$) erf\"ullt jedoch $C_1 \wedge C_2$. Damit folgt
der Doppelresolvent nicht logisch aus den Pr\"amissen, die Regel ist
\textbf{inkorrekt}. \hfill$\square$%
}}

\newpage

%% ============================================================
%% AUFGABE 2
%% ============================================================
\subsection*{Aufgabe 2 \textnormal{(3 Punkte)} -- Resolution \& DPLL f\"ur $F$}

\textbf{Was ist eine CNF?}
CNF hei\ss t \emph{konjunktive Normalform}: eine Formel, die ein gro\ss es
UND von Klammern ist, wobei jede Klammer eine Oder-Verkn\"upfung
(eine Klausel) ist. Eine CNF ist genau dann wahr, wenn jede einzelne
Klausel wahr ist. F\"ur DPLL und f\"ur Resolution muss die Formel zuerst in
CNF gebracht werden.

\medskip
\textbf{Resolution als Beweisverfahren.}
Die Idee: Aus der Klauselmenge werden so lange neue Klauseln durch
Resolution abgeleitet, bis entweder die leere Klausel $\emptyset$
auftaucht (dann ist die Formel unerf\"ullbar) oder keine neue Klausel
mehr entsteht (dann ist sie erf\"ullbar). Ziel ist hier also, einen
Ableitungsweg von den Ausgangsklauseln bis zu $\emptyset$ zu finden.

\medskip
\textbf{DPLL kurz.}
DPLL setzt Variablen Schritt f\"ur Schritt auf $0$ oder $1$ und
vereinfacht danach die Klauseln:
\begin{itemize}[nosep]
\item Eine Klausel, die wahr wird, f\"allt weg.
\item Ein Literal, das falsch wird, wird nur aus seiner Klausel entfernt.
\item \emph{Unit Propagation}: Bleibt in einer Klausel nur ein Literal
      \"ubrig (eine Unit-Klausel), muss dieses Literal wahr werden.
\item Entsteht eine leere Klausel $\emptyset$, scheitert der Zweig. DPLL
      geht zum letzten Verzweigungspunkt zur\"uck und probiert den
      anderen Wert.
\end{itemize}

\medskip
\textbf{Aufgabenstellung.}
Entscheide mittels Resolution, ob die Formel
\[
F \;\equiv\; \neg(X \to \neg Y)
\;\wedge\; \neg(X \wedge Y \wedge \neg A)
\;\wedge\; (X \wedge A \to B)
\;\wedge\; \neg B
\]
erf\"ullbar ist. Zus\"atzlich sind die Schritte anzugeben, die DPLL
machen w\"urde.

\bigskip
\textbf{Schritt 1: CNF bestimmen.}

Bei der Umformung werden zwei Regeln immer wieder gebraucht:
\begin{itemize}[nosep]
\item Implikation aufl\"osen: $P \to Q \;\equiv\; \neg P \vee Q$.
\item De Morgan: $\neg(A \vee B) \equiv \neg A \wedge \neg B$ und
      $\neg(A \wedge B) \equiv \neg A \vee \neg B$.
\end{itemize}
Doppelte Negation $\neg\neg A \equiv A$ wird stillschweigend mit verwendet.
Die vier Konjunkte werden einzeln umgeformt.

\medskip
\emph{Erstes Konjunkt: $\neg(X \to \neg Y)$.}
Erst die Implikation aufl\"osen:
\begin{align*}
X \to \neg Y &\equiv \neg X \vee \neg Y.
\end{align*}
Damit:
\begin{align*}
\neg(X \to \neg Y)
&\equiv \neg(\neg X \vee \neg Y)
&& \text{(Implikation eingesetzt)} \\
&\equiv \neg(\neg X) \wedge \neg(\neg Y)
&& \text{(De Morgan)} \\
&\equiv X \wedge Y.
&& \text{(doppelte Negation)}
\end{align*}
Das ergibt zwei Unit-Klauseln $\{X\}$ und $\{Y\}$.

\medskip
\emph{Zweites Konjunkt: $\neg(X \wedge Y \wedge \neg A)$.}
De Morgan f\"ur eine Drei-Konjunktion:
\begin{align*}
\neg(X \wedge Y \wedge \neg A)
&\equiv \neg X \vee \neg Y \vee \neg(\neg A)
&& \text{(De Morgan)} \\
&\equiv \neg X \vee \neg Y \vee A.
&& \text{(doppelte Negation)}
\end{align*}
Klausel: $\{\neg X, \neg Y, A\}$.

\medskip
\emph{Drittes Konjunkt: $X \wedge A \to B$.}
Implikation aufl\"osen, dann De Morgan auf $\neg(X \wedge A)$:
\begin{align*}
X \wedge A \to B
&\equiv \neg(X \wedge A) \vee B
&& \text{(Implikation aufgel\"ost)} \\
&\equiv (\neg X \vee \neg A) \vee B
&& \text{(De Morgan)} \\
&\equiv \neg X \vee \neg A \vee B.
\end{align*}
Klausel: $\{\neg X, \neg A, B\}$.

\medskip
\emph{Viertes Konjunkt: $\neg B$.}
Bereits eine Klausel: $\{\neg B\}$.

\medskip
\textbf{Klauselmenge.}
\[
\begin{array}{rclcl}
C_1 &\equiv& X                          &=& \{X\},\\
C_2 &\equiv& Y                          &=& \{Y\},\\
C_3 &\equiv& \neg X \vee \neg Y \vee A  &=& \{\neg X, \neg Y, A\},\\
C_4 &\equiv& \neg X \vee \neg A \vee B  &=& \{\neg X, \neg A, B\},\\
C_5 &\equiv& \neg B                     &=& \{\neg B\}.
\end{array}
\]

\bigskip
\textbf{Schritt 2: Resolution.}

Es gibt drei Unit-Klauseln ($C_1, C_2, C_5$), die jeweils ein einzelnes
Literal erzwingen. Die Strategie ist daher, diese Unit-Klauseln nach und
nach mit den l\"angeren Klauseln $C_3$ und $C_4$ zu resolvieren und so die
l\"angeren Klauseln zu k\"urzen, bis am Ende zwei widerspr\"uchliche
Unit-Klauseln \"ubrig sind, deren Resolvent $\emptyset$ ist.

\begin{center}
\small
\renewcommand{\arraystretch}{1.3}
\begin{tabular}{p{0.08\textwidth}p{0.27\textwidth}p{0.57\textwidth}}
\toprule
Schritt & Resolvent & Wie der Resolvent entsteht \\
\midrule
(a)
& $C_6 = \{\neg Y, A\}$
& Resolution von $C_1 = \{X\}$ und $C_3 = \{\neg X, \neg Y, A\}$ \"uber $X$.\\
&& Aus $C_1$ verschwindet $X$, aus $C_3$ verschwindet $\neg X$.\\
&& Resolvent: $\emptyset \cup \{\neg Y, A\} = \{\neg Y, A\}$. \\
\midrule
(b)
& $C_7 = \{A\}$
& Resolution von $C_2 = \{Y\}$ und $C_6 = \{\neg Y, A\}$ \"uber $Y$.\\
&& Aus $C_2$ verschwindet $Y$, aus $C_6$ verschwindet $\neg Y$.\\
&& Resolvent: $\emptyset \cup \{A\} = \{A\}$. \\
\midrule
(c)
& $C_8 = \{\neg X, \neg A\}$
& Resolution von $C_4 = \{\neg X, \neg A, B\}$ und $C_5 = \{\neg B\}$ \"uber $B$.\\
&& Aus $C_4$ verschwindet $B$, aus $C_5$ verschwindet $\neg B$.\\
&& Resolvent: $\{\neg X, \neg A\} \cup \emptyset = \{\neg X, \neg A\}$. \\
\midrule
(d)
& $C_9 = \{\neg A\}$
& Resolution von $C_1 = \{X\}$ und $C_8 = \{\neg X, \neg A\}$ \"uber $X$.\\
&& Aus $C_1$ verschwindet $X$, aus $C_8$ verschwindet $\neg X$.\\
&& Resolvent: $\emptyset \cup \{\neg A\} = \{\neg A\}$. \\
\midrule
(e)
& $\emptyset$
& Resolution von $C_7 = \{A\}$ und $C_9 = \{\neg A\}$ \"uber $A$.\\
&& Aus $C_7$ verschwindet $A$, aus $C_9$ verschwindet $\neg A$.\\
&& Resolvent: $\emptyset \cup \emptyset = \emptyset$ -- die leere Klausel.\\
\bottomrule
\end{tabular}
\end{center}

Die leere Klausel ist erreicht, also ist $F$ \textbf{unerf\"ullbar}.

\bigskip
\textbf{Schritt 3: DPLL.}

Es gibt am Anfang gleich drei Unit-Klauseln: $\{X\}$, $\{Y\}$, $\{\neg B\}$.
Die werden nacheinander durch Unit Propagation abgearbeitet, ohne dass
verzweigt werden muss.

\begin{center}
\small
\renewcommand{\arraystretch}{1.3}
\begin{tabular}{p{0.18\textwidth}p{0.22\textwidth}p{0.52\textwidth}}
\toprule
Schritt & Belegung & Was mit den Klauseln passiert \\
\midrule
Start
& noch keine Werte
& Klauselmenge: $\{X\},\;\{Y\},\;\{\neg X, \neg Y, A\},\;\{\neg X, \neg A, B\},\;\{\neg B\}$. \\
\midrule
Unit $\{X\}$
& $X = 1$
& $\{X\}$ wird wahr und f\"allt weg.\\
&& In $\{\neg X, \neg Y, A\}$ ist $\neg X$ falsch und wird entfernt -- es bleibt $\{\neg Y, A\}$.\\
&& In $\{\neg X, \neg A, B\}$ ist $\neg X$ falsch und wird entfernt -- es bleibt $\{\neg A, B\}$.\\
&& Rest: $\{Y\},\;\{\neg Y, A\},\;\{\neg A, B\},\;\{\neg B\}$. \\
\midrule
Unit $\{Y\}$
& $X = 1,\; Y = 1$
& $\{Y\}$ wird wahr und f\"allt weg.\\
&& In $\{\neg Y, A\}$ ist $\neg Y$ falsch und wird entfernt -- es bleibt $\{A\}$.\\
&& Rest: $\{A\},\;\{\neg A, B\},\;\{\neg B\}$. \\
\midrule
Unit $\{A\}$
& $X = 1,\; Y = 1,\; A = 1$
& $\{A\}$ wird wahr und f\"allt weg.\\
&& In $\{\neg A, B\}$ ist $\neg A$ falsch und wird entfernt -- es bleibt $\{B\}$.\\
&& Rest: $\{B\},\;\{\neg B\}$. \\
\midrule
Unit $\{B\}$
& $X = 1,\; Y = 1,\; A = 1,\; B = 1$
& $\{B\}$ wird wahr und f\"allt weg.\\
&& In $\{\neg B\}$ wird $\neg B$ falsch und entfernt -- es bleibt die leere Klausel $\emptyset$. Konflikt. \\
\bottomrule
\end{tabular}
\end{center}

Da alle Belegungen durch Unit Propagation \emph{erzwungen} waren, gibt es
keinen Verzweigungspunkt, an dem DPLL einen anderen Wert probieren
k\"onnte. Das Backtracking l\"auft ins Leere -- DPLL meldet
\emph{unerf\"ullbar}.

\medskip
\noindent\fbox{\parbox{\dimexpr\textwidth-2\fboxsep-2\fboxrule}{%
\textbf{Ergebnis Aufgabe 2.}\\[2pt]
Resolution leitet die leere Klausel $\emptyset$ ab; DPLL l\"auft durch
reine Unit Propagation in einen Konflikt, ohne dass verzweigt werden
muss. Beide Verfahren liefern dasselbe Resultat: die Formel
\[
F \;\equiv\; \neg(X \to \neg Y)
\wedge \neg(X \wedge Y \wedge \neg A)
\wedge (X \wedge A \to B)
\wedge \neg B
\]
ist \textbf{unerf\"ullbar}.%
}}

\newpage

%% ============================================================
%% AUFGABE 3
%% ============================================================
\subsection*{Aufgabe 3 \textnormal{(3 Punkte)} -- Resolution \& DPLL f\"ur $G$}

\textbf{Vorgehen.} Wie in Aufgabe 2: CNF, dann Resolution + DPLL. Pointe:
in der Klauselmenge kommt $\neg B$ nirgends vor.

\medskip
\textbf{Aufgabenstellung.}
Entscheide mittels Resolution, ob
\[
G \;\equiv\; (X \vee Y)
\;\wedge\; \neg(X \wedge Y \wedge \neg A)
\;\wedge\; (X \wedge A \to B)
\;\wedge\; (A \vee B)
\]
erf\"ullbar ist. Zus\"atzlich die Schritte angeben, die DPLL machen w\"urde.

\bigskip
\textbf{Schritt 1: CNF.}
Jeder Konjunkt wird mit den bekannten Regeln (De Morgan, $P \to Q \equiv \neg P \vee Q$)
direkt zu seiner Klausel:

\[
G \;\equiv\;
\underbrace{(X \vee Y)}_{\textstyle C_1 = \{X,\,Y\}}
\,\wedge\,
\underbrace{\neg(X \wedge Y \wedge \neg A)}_{\textstyle C_2 = \{\neg X,\,\neg Y,\,A\}}
\,\wedge\,
\underbrace{(X \wedge A \to B)}_{\textstyle C_3 = \{\neg X,\,\neg A,\,B\}}
\,\wedge\,
\underbrace{(A \vee B)}_{\textstyle C_4 = \{A,\,B\}}.
\]

\bigskip
\textbf{Schritt 2: Resolution.}
Vorkommen jedes Literals in der Klauselmenge:

\begin{center}
\renewcommand{\arraystretch}{1.25}
\begin{tabular}{r|cc|cc|cc|cc}
\toprule
                              & $X$        & $\neg X$   & $Y$        & $\neg Y$   & $A$        & $\neg A$   & $B$        & $\neg B$ \\
\midrule
$C_1 = \{X, Y\}$              & \checkmark &            & \checkmark &            &            &            &            &          \\
$C_2 = \{\neg X, \neg Y, A\}$ &            & \checkmark &            & \checkmark & \checkmark &            &            &          \\
$C_3 = \{\neg X, \neg A, B\}$ &            & \checkmark &            &            &            & \checkmark & \checkmark &          \\
$C_4 = \{A, B\}$              &            &            &            &            & \checkmark &            & \checkmark &          \\
\bottomrule
\end{tabular}
\end{center}

Spalte $\neg B$ ist leer:
\[
\neg B \text{ kommt nirgends vor}
\;\;\Longrightarrow\;\;
B \text{ ist nicht eliminierbar}
\;\;\Longrightarrow\;\;
\emptyset \text{ nicht ableitbar}
\;\;\Longrightarrow\;\;
G \text{ \textbf{erf\"ullbar}}.
\]

\bigskip
\textbf{Schritt 3: DPLL.}
Keine Unit-Klausel am Anfang, also Verzweigung. Pfad $X = 1 \to A = 1 \to B = 1$:

\begin{center}
\renewcommand{\arraystretch}{1.35}
\begin{tabular}{ll}
\toprule
\textbf{Schritt} & \textbf{Klauselmenge nach Vereinfachung} \\
\midrule
Start              & $\{X, Y\},\;\{\neg X, \neg Y, A\},\;\{\neg X, \neg A, B\},\;\{A, B\}$ \\
Wahl $X = 1$       & $\{\neg Y, A\},\;\{\neg A, B\},\;\{A, B\}$ \\
Wahl $A = 1$       & $\{B\}$ \\
Unit $B = 1$       & alles erf\"ullt \\
\bottomrule
\end{tabular}
\end{center}

\textbf{Modell:} $X = 1,\; A = 1,\; B = 1$, $Y$ frei.

\medskip
\noindent\fbox{\parbox{\dimexpr\textwidth-2\fboxsep-2\fboxrule}{%
\textbf{Ergebnis Aufgabe 3.}\\[2pt]
$\neg B$ kommt in keiner Klausel vor $\Rightarrow$ Resolution kann
$\emptyset$ nicht ableiten $\Rightarrow$ $G$ \textbf{erf\"ullbar}.
DPLL-Modell: $X = 1,\; A = 1,\; B = 1$, $Y$ beliebig.%
}}

\end{document}
